\scanpage{33}
(quitte à faire une extension finie sur $S$ et remplacer $G$ par $\overline{G_{\deleted{\unsure{\mathrm{réd}}}1\,\mathrm{réd}}^{\circ}}$] que $G_1$ est \inserted{connexe et} lisse $/k(s_1)$, enfin [quitte à remplacer $G$ par $G$ $-$ composantes non neutres de $G_0$] que $G_0$ est également connexe [mais peut-être pas lisse].

a) L'inégalité $d\geq d'$ bien connue, mise pour mémoire, elle se réduit \deleted{\unsure{alors}}\inserted{ici} (cas plat) à l'égalité $d=d'$. L'inégalité $\rho_{aff}\geq\rho'_{aff}$ \uline{équivaut alors} à $\rho_{ab}\leq\rho'_{ab}$ [\unsure{car|or} $\rho_{aff}+\rho_{ab}=\rho'_{aff}+\rho'_{ab}=d$], que nous allons traiter plus bas.

b) \deleted{Considérons} On peut supposer que $G_0$ admet un tore maximal splitté $T_0$ (car on peut supposer $k(s_0)$ alg. clos]. \deleted{\unsure{les|Les} $ {}_nT_0$ Soit} $C_0$ le Cartan associé. \deleted{\unsure{si|on peut} \unsure{premier|prolonge|prolonger} \unsure{premier} à \unsure{la caract.|la car.} résiduelle.} Pour tt $n$ premier

\margin{\deleted{$G$ lisse$/S$}}
\deleted{Alors pour les puiss. \unsure{$n\ell$|$\ell$|de $\ell$} de $n$,}
à la caractéristique résiduelle, considérons

\deleted{$ {}_n C_0=\operatorname{Ker}$ [\unsure{mult$^n$|$n^{\mathrm{ième}}$} de $C_0$]. \unsure{Comme} $C_0/T_0=\deleted{\unsure{G}}M_0$ est une extension d'une VA $A_0$ par un groupe unipotent \inserted{lisse} connexe $U_0$, on voit que les extensions \unsure{correspondantes} de $ {}_n A_0$ par $U_0$ splittent \uuline{canoniquement}, \deleted{\unsure{d'où|donc} \unsure{son extension|des extensions}}\inserted{l'image étant nlle, d'où}}
% archive sheet number: 32
